% 08_biophysics.tex
% Ch.8: Protein backbone impedance

\chapter{Biophysics: Protein Folding}
\label{ch:biophysics}


\begin{objectivebox}
\begin{itemize}
    \item Derive the C$\alpha$--C$\alpha$ virtual bond distance $d_0 \approx 3.80$~\AA\ from the AVE lattice pitch $\ell_{\text{node}}$ through a chain of topological identities involving the Bohr radius and covalent bond lengths.
    \item Classify the 20 canonical amino acids by their sidechain impedance relative to the backbone baseline ($Z_{\text{backbone}} \approx 7$).
    \item Interpret topological structural evaluations as a Level~6 Emergence validation using only \texttt{impedance()}, \texttt{saturation\_factor()}, and \texttt{reflection\_coefficient()}.
    \item Evaluate the Molecular Chiral FRET Parallax as a future falsification target, understanding why terrestrial baselines are insufficient.
\end{itemize}
\end{objectivebox}

\section{Protein Backbone: From Proton Radius to Folding}
\label{sec:protein_bridge}

The protein $S_{11}$ topological derivation uses the C$\alpha$--C$\alpha$
virtual bond distance $d_0 \approx 3.80$~\AA\ as its fundamental
length scale.  This traces back to the AVE axioms via:
\begin{resultbox}{Proton-to-Backbone Length Scale}
\begin{equation}
    d_p = \frac{4\hbar}{m_p c} \approx 0.841\;\text{fm}
    \;\xrightarrow{\text{nuclear bond}}
    r_{\text{cov}}
    \;\xrightarrow{\text{backbone}}
    d_0 \approx 3.80\;\text{\AA}
\end{equation}
\end{resultbox}
The ratio $d_0 / a_0 \approx 7.18$ (where $a_0 = \ell_{\text{node}}
/\alpha$ is the Bohr radius) suggests a geometric origin in the
peptide repeat unit.  The amino acid impedance table $Z_{\text{topo}}$
assigns complex impedances (R + jX) to each of the 20 residues,
with $Q_{\text{backbone}} \approx 7$ suppressing reactive coupling.

\begin{figure}[H]
    \centering
    \includegraphics[width=0.95\textwidth]{protein_backbone_impedance.pdf}
    \caption{Protein backbone impedance from AVE axioms.
    \textbf{Left:} Full derivation chain from $\ell_{node}$ to
    the C$\alpha$--C$\alpha$ peptide bond length $d_0 = 3.80$~\AA.
    \textbf{Centre:} Amino acid impedance spectrum --- hydrophobic
    (gold), positive/inductive (blue), negative/capacitive (red).
    \textbf{Right:} Local inter-residue reflection coefficient
    $\Gamma$ along a sample peptide sequence; high $|\Gamma|$
    indicates folding-driving impedance mismatch.}
    \label{fig:protein_backbone}
\end{figure}

\section{Derivation Chain: From Lattice Pitch to Bond Length}

The C$\alpha$--C$\alpha$ virtual bond distance $d_0 \approx 3.80$~\AA\ is derived through a chain of AVE identities:

\begin{enumerate}
    \item \textbf{Lattice pitch:} $\ell_{node} = \hbar / (m_e c) \approx 3.862 \times 10^{-13}$ m.
    \item \textbf{Bohr radius:} $a_0 = \ell_{node} / \alpha \approx 5.292 \times 10^{-11}$ m (the characteristic extension of the unknot electron).
    \item \textbf{Covalent bond length:} C--N and C$\alpha$--C bonds are set by the overlap of atomic ground states at the equilibrium impedance-matching distance: $r_{cov} \approx 1.47$~\AA\ (C--N), 1.52~\AA\ (C--C), 1.33~\AA\ (C=O).
    \item \textbf{Backbone repeat:} The peptide unit (N--C$\alpha$--C--N) has a fixed geometry from the $sp^3$ tetrahedral bond angle ($\theta = 109.47^\circ$) and the planar amide constraint ($\omega = 180^\circ$). The resulting C$\alpha$--C$\alpha$ distance:
    \begin{equation}
        d_0 = \sqrt{r_{C\alpha C}^2 + r_{CN}^2 + 2\, r_{C\alpha C}\, r_{CN} \cos(180^\circ - \theta)} \approx 3.80 \text{ \AA}
    \end{equation}
\end{enumerate}

Every step uses constants from the fundamental AVE topological derivations (Volume~I) or standard NERF rigid-body geometry. No free parameter is introduced at the biological scale.

\section{Amino Acid Impedance Classification}

The 20 canonical amino acids are classified by their sidechain impedance relative to the backbone baseline ($Z_{backbone} \approx 7$):

\begin{itemize}
    \item \textbf{Hydrophobic (low $Z$):} Ala, Val, Leu, Ile, Pro, Phe, Trp, Met --- compact sidechains with low reactive coupling. $|\Gamma| \approx 0$ (impedance-matched, minimal folding drive).
    \item \textbf{Polar/Charged (high $|Z|$):} Glu, Asp (capacitive), Lys, Arg (inductive) --- large $|\Gamma|$, strong folding drive at sequence positions where impedance mismatch is maximal.
    \item \textbf{Special:} Gly ($Z \approx 0$, maximum flexibility), Cys (disulphide bridge, reactive short-circuit between distant residues).
\end{itemize}

\subsection{RF-to-Biology Translation Dictionary}
\label{sec:rf_biology_dictionary}

The protein folding derivation uses radio-frequency (RF) circuit terminology throughout. Table~\ref{tab:rf_bio_dict} provides a persistent reference for readers crossing between electrical engineering and biochemistry.

\begin{table}[h]
\centering
\renewcommand{\arraystretch}{1.3}
\begin{tabularx}{\textwidth}{@{}llX@{}}
\toprule
\textbf{EE / RF Term} & \textbf{Biochemistry Term} & \textbf{AVE Operator} \\
\midrule
Transmission line      & Peptide backbone           & $Z_c = \sqrt{L/C}$ (Op1) \\
Characteristic impedance & Backbone stiffness       & $Z_\text{topo}$ per residue \\
Shunt stub             & Side chain (R-group)       & Reactive filter load \\
S$_{11}$ (reflection)  & Folding stability          & Op3: $\Gamma = (Z_2-Z_1)/(Z_2+Z_1)$ \\
Impedance matching     & Native fold geometry       & $\lambda_\text{min}(S^\dagger S) \to 0$ (Op6) \\
Bandpass filter        & Sequence selectivity       & Cascaded ABCD matrices \\
Q factor               & Conformational rigidity    & $Q = \omega_0 / \Delta\omega$ \\
VSWR                   & Folding frustration        & $(1+|\Gamma|)/(1-|\Gamma|)$ \\
Mutual inductance      & H-bond coupling            & $K_{HB} = \Gamma^2 \alpha\hbar c$ \\
Nodal admittance $[Y]$ & Residue connectivity       & Op5: $[S] = ([Y]+I)^{-1}([Y]-I)$ \\
Resistive loss         & Solvent damping            & $G_\text{solvent} = \gamma \xi^2$ \\
Noise floor            & Thermal cytosol bath       & Wien $kT$ at 310 K \\
\bottomrule
\end{tabularx}
\caption{RF-to-Biology translation dictionary. Every protein folding term maps directly to a standard electrical engineering concept via one of the 15 universal AVE operators.}
\label{tab:rf_bio_dict}
\end{table}

\section{Biophysical Validation Bounds}

To validate the strict limits of macroscopic biophysics, topological bounds must predictably cascade from an extended chain configuration to a rigid $\beta$-hairpin equilibrium:
\begin{resultbox}{Biophysical Falsification}
\begin{equation}
    \text{RMSD}_{\text{backbone}} = 2.59 \text{ \AA} \quad \text{(sub-3 \AA: Level 6 Emergence)}
\end{equation}
\end{resultbox}

This validation uses the \textit{same three functions} as every other chapter in this volume: \texttt{impedance()}, \texttt{saturation\_factor()}, and \texttt{reflection\_coefficient()}. The protein backbone is not ``biology''---it is a transmission line.

\section{Molecular Chiral FRET Parallax}
To experimentally map the continuous boundary of Axiom 2 (Chirality) into living biophysics, the AVE framework proposes the \textbf{Molecular Chiral FRET Parallax} test. Standard biochemistry assumes the left-handed bias of terrestrial amino acids is a random, isolated historical accident. AVE dictates it is the macroscopic manifestation of the $h_{\perp}$ chiral metric bias structurally enforcing the Ramachandran exclusion zones.

By attaching Cy3/Cy5 fluorophore pairs to the terminal ends of a rigid structural $\alpha$-helix and measuring the F\"orster Resonance Energy Transfer (FRET) efficiency, researchers can map the rigid spatial bounds of the $sp^3$ tetrahedral carbon limits. Under AVE axioms, transitioning from Earth's $1G$ gravitational VSWR ($\varepsilon_{11} \approx 7 \times 10^{-10}$) to microgravity relaxes the background metric strain by $\Delta r / r = \alpha \cdot \varepsilon_{11} \approx 5 \times 10^{-12}$. For a $5\text{ nm}$ fluorophore baseline, this corresponds to a sub-attometer ($\sim 10^{-20}\text{ m}$) physical relaxation---far below any current single-molecule detection threshold.

\textbf{Current Status:} The Molecular Chiral FRET Parallax is \textit{currently unfalsifiable} at terrestrial gravitational baselines. The honest AVE prediction confirms the chiral metric bias is physically real but requires either (a) a compact-object environment (e.g., Sirius B, $\varepsilon_{11} \sim 10^{-4}$) or (b) a coherent resonant amplification mechanism to reach detectable signal levels. This test is catalogued as a \textit{future falsification target} pending access to extreme gravitational environments.

\begin{summarybox}
\begin{itemize}
    \item The protein backbone is derived from the proton radius through a chain of AVE identities: $\ell_{\text{node}} \to a_0 \to r_{\text{cov}} \to d_0 \approx 3.80$~\AA, introducing zero free parameters.
    \item The 20 canonical amino acids are classified by sidechain impedance relative to the backbone baseline ($Z_{\text{backbone}} \approx 7$), mapping hydrophobic, polar, and charged residues to reactive coupling strengths.
    \item Macroscopic rigid proteins establish theoretical bounds to sub-3~\AA\ precision using only \texttt{impedance()}, \texttt{saturation\_factor()}, and \texttt{reflection\_coefficient()}---the same three functions used at every other scale.
\end{itemize}
\end{summarybox}

\begin{exercisebox}
\begin{enumerate}
    \item Derive the C$\alpha$--C$\alpha$ virtual bond distance $d_0$ from the peptide unit bond lengths and $sp^3$ tetrahedral angle. Verify that no free parameter is introduced at the biological scale.
    \item Explain why the Molecular Chiral FRET Parallax test is currently unfalsifiable, and estimate the minimum gravitational VSWR required for a detectable signal.
\end{enumerate}
\end{exercisebox}
